\chapter{The Interactive Implementation}\label{ch:lisp}

\index{interactive implementation}
The interactive implementation runs Dylan inside a Common Lisp system,
SBCL. You type a Dylan expression, it is compiled to Common Lisp at once,
the Lisp compiles that to machine code and runs it, and the value is
printed. Nothing about the program is fixed in advance: classes can gain
slots, generic functions can gain and lose methods, and definitions can
be replaced while the program runs, all as the 1992 manual describes.
That makes it the natural place to learn the language, to try out an
idea, and to follow the DIRM's examples.

This chapter describes how the implementation is built, how to load and
use it, how Dylan values look to Lisp and how to call Lisp from Dylan,
how it is made fast, and how it differs from the compiler to C.

\section{How it works}\label{sec:lisp-architecture}

\index{interactive implementation!architecture}
The implementation began in 1992 as Thomas, a Dylan system written in
Scheme at Digital Equipment Corporation's Cambridge Research Laboratory.
This project translated it into Common Lisp and has since corrected it
against the manual and rebuilt much of it for speed. Every source file in
\file{src/} still names the Thomas file it grew from.

An expression goes through four stages.

\begin{enumerate}
\item \textbf{Reading.} Common Lisp's reader cannot read Dylan: a colon
  is a package marker to it, and it knows nothing of \code{\#t} and
  \code{\#f}. A small reader of Dylan's own (in \file{compat.lisp}) turns
  text into lists, symbols, numbers and strings. A Dylan symbol becomes a
  Lisp symbol in the \code{DYLAN} package; a keyword \code{foo:} becomes
  the Lisp keyword \code{:FOO}.
\item \textbf{Compiling to Lisp.} The compiler (\file{compiler.lisp}
  and the \file{comp-*.lisp} files) turns each Dylan form into a Lisp
  form: \code{begin} into \code{progn}, \code{set!} into \code{setq}, a
  method into a \code{lambda}, a call into a call through the runtime's
  dispatcher. Each module variable becomes a Lisp special variable
  with a name of its own, \code{!MODULE/X/} for a variable \code{x} of
  \code{dylan-user} and \code{!MODULE/X@M/} for one owned by module
  \code{m}.
\item \textbf{Compiling to machine code.} The Lisp form is handed to
  SBCL's own compiler, which turns it into native code.
\item \textbf{Running.} The code runs against the runtime library, the
  \file{runtime-*.lisp} files, which implement classes, generic
  dispatch, collections, numbers and conditions.
\end{enumerate}

\noindent
The files of \file{src/}, in the order they are loaded, do these jobs:

\begin{longtable}{@{}>{\raggedright\arraybackslash}p{0.34\textwidth}>{\raggedright\arraybackslash}p{0.6\textwidth}@{}}
\toprule
File & Role \\
\midrule
\endhead
\file{package.lisp} & the \code{DYLAN} package and its exported entry points \\
\file{globals.lisp} & declares the runtime's global variables, which nearly every file refers to before the file that defines them is loaded \\
\file{compat.lisp} & the Dylan reader and printer, weak tables, escapes, sorting and the interface to Common Lisp's condition system \\
\file{support.lisp} & small utilities shared by the compiler and runtime \\
\file{comp-util.lisp} & the compiler's utilities: names, variables, parameter lists \\
\file{comp-module.lisp} & modules, \code{define-module} and the \code{module:} file header (Chapter~\ref{ch:modules}) \\
\file{compiler.lisp} & the heart of the compiler, the table of special forms, the predefined names, \code{compile-file} and \code{compile-to-file} \\
\file{comp-method.lisp} & \code{method}, \code{define-method}, \code{define-generic-function} \\
\file{comp-class.lisp} & \code{define-class} and \code{define-slot} \\
\file{comp-sf.lisp} & the other special forms: \code{bind}, \code{for}, \code{select} and the rest \\
\file{comp-exc.lisp} & \code{handler-bind}, \code{handler-case}, \code{bind-exit}, \code{unwind-protect} \\
\file{class.lisp} & classes, slots, instances and \code{make} \\
\file{generic.lisp} & generic functions and method dispatch, with its caches \\
\file{class-structure.lisp} & the predefined class heterarchy \\
\file{runtime-top.lisp}, \file{runtime-internal.lisp} & common parameter lists and the procedures compiled code calls \\
\file{runtime-methods.lisp}, \file{runtime-functions.lisp} & the predefined methods and functions \\
\file{runtime-extensions.lisp} & \code{lisp-variable}, \code{lisp-procedure} \\
\file{runtime-bitstrings.lisp} & \code{logand} and the other bit operations \\
\file{runtime-collections*.lisp} & collections: the iteration protocol, the generic operations, and one file for each kind (lists, vectors, strings, deques, ranges, tables, arrays) \\
\file{runtime-exceptions.lisp} & conditions, handlers and restarts \\
\file{runtime-printing.lisp} & streams, \code{print-object}, \code{format} and the other printing functions (Chapter~\ref{ch:printing}) \\
\file{runtime-extended-float.lisp} & \code{<extended-float>} arithmetic (Chapter~\ref{ch:numbers}) \\
\file{runtime-fast-paths.lisp} & turns on the guarded fast paths (Section~\ref{sec:lisp-performance}) \\
\file{rep.lisp} & the read-eval-print loop \\
\file{load-dylan.lisp}, \file{load-compiler.lisp}, \file{load-runtime.lisp} & loaders \\
\file{build-dylan.lisp} & saves the executable \file{bin/dylan} \\
\file{full-test.lisp} & the regression test driver \\
\bottomrule
\end{longtable}

\noindent
\file{src/README} describes the translation from Scheme in more detail.

\section{Loading it}\label{sec:lisp-loading}

\index{dylan@\cmd{dylan}}\index{build-dylan.lisp@\file{build-dylan.lisp}}
The usual way in is the executable \file{bin/dylan}. Running \cmd{make}
in the repository's top directory builds it: \file{src/build-dylan.lisp}
loads the whole system into SBCL and saves it, with SBCL, as one
executable that starts in the read-eval-print loop.

\begin{shell}
$ make
$ bin/dylan

Entering Dylan read-eval-print-loop.
Leave by typing "(exit)"

?
\end{shell}

\noindent
The executable needs nothing else from the repository, so it may be
copied or linked anywhere. It does not see later changes to \file{src/};
run \cmd{make} again to rebuild it. On Windows, without \cmd{make}, run
\cmd{sbcl -{}-script src/build-dylan.lisp}, which writes
\file{bin\textbackslash dylan.exe}.

\index{load-dylan.lisp@\file{load-dylan.lisp}}
The other way is to start a plain SBCL in the repository's top directory
and load the system into it, which is how \cmd{make} itself begins:

\begin{shell}
$ sbcl
* (load "src/load-dylan.lisp")
;Dylan loaded.  (dylan:rep) starts a read-eval-print loop.
\end{shell}

\index{fasl cache}
The first time, this compiles every source file into a \term{fasl}
(Lisp's compiled-file format), which takes a little while. The fasls are
kept in \file{src/fasl/\meta{lisp}-\meta{version}/}, for example
\file{src/fasl/SBCL-2-6-3/}, one directory for each Lisp and version
since their fasls are incompatible. Later loads use them directly. When a
source file changes, it is recompiled together with every file loaded
after it, since those may depend on its macros and inline functions.

Three loaders exist, for the three ways of using the system that Thomas
offered:

\begin{description}
\item[\file{load-dylan.lisp}] the interactive environment: the
  read-eval-print loop.
\item[\file{load-compiler.lisp}] the compiler to Lisp:
  \code{dylan:compile-file} and \code{dylan:compile-to-file}.
\item[\file{load-runtime.lisp}] what compiled programs need to run.
\end{description}

\noindent
In this version the three load the same thing, the whole system;
\file{load-compiler.lisp} and \file{load-runtime.lisp} simply load
\file{load-dylan.lisp}. Use whichever names your purpose. A loader can
be given on the command line as well:

\begin{shell}
$ sbcl --eval '(load "src/load-dylan.lisp")' --eval '(dylan:rep)'
\end{shell}

\section{The read-eval-print loop}\label{sec:lisp-rep}

\index{read-eval-print loop}\index{rep@\code{dylan:rep}}
\file{bin/dylan} starts the loop at once; in a Lisp with the system
loaded, \code{(dylan:rep)} starts it. It prompts with \code{?}, reads one
Dylan form, compiles and runs it, and prints the value after
\code{;Value:}. A form may span several lines; the loop reads until the
form is complete.

\begin{repl}
* (dylan:rep)

Entering Dylan read-eval-print-loop.
Leave by typing "(exit)"

? (define-method double ((n <number>)) (+ n n))

;Value: double

? (double 21)

;Value: 42

? (values 1 2 3)

;Value[1]: 1
;Value[2]: 2
;Value[3]: 3

? (values)


? dylan:done
DYLAN::|DYLAN:DONE|
*
\end{repl}

\noindent
A form that returns several values shows each on its own
\code{;Value[\meta{n}]:} line; one that returns none shows nothing.
Values are printed as \code{print} prints them (Chapter~\ref{ch:printing}),
and output the form itself writes appears before the \code{;Value:} line.
A definition answers the name it defined.

Typing \code{dylan:done}, or ending the input (Ctrl-D), leaves the
loop. In \file{bin/dylan} that ends the program. In a Lisp you loaded
the system into, it returns to Lisp's \code{*} prompt; the Lisp prints
the symbol the loop returned. Everything the loop defined remains: call
\code{(dylan:rep)} again and carry on where you left off. \code{dylan:done}
may also be typed in parentheses, as \code{(dylan:done)}.

\index{exit@\code{(exit)}}
Typing \code{(exit)} quits altogether: it ends \file{bin/dylan}, and in a
Lisp you loaded the system into, it ends the Lisp as well. It is a command
of the loop, like \code{dylan:done}, not a Dylan function; a variable
named \code{exit} is still evaluated when typed bare. Typing \code{(dylan:in-module \meta{name})} at the
\code{?} prompt changes the module later input is compiled in
(Section~\ref{sec:modules-lisp}).

Redefining things is normal at the loop. A second \code{define} of a
variable replaces its value with a warning, a second \code{define-method}
with the same specializers replaces the method, and a second
\code{define-class} redefines the class in place:

\begin{repl}
? (define v 1)
;Value: v

? (define v 2)
WARNING: redefining variable v
;Value: v

? (define-method describe ((x <integer>)) "an integer")
;Value: describe

? (define-method describe ((x <integer>)) "a whole number")
;Value: describe

? (describe 3)
;Value: "a whole number"
\end{repl}

\index{empty-environment"!@\code{dylan:empty-environment"!}}
From Lisp, \code{(dylan:empty-environment!)} makes the loop forget the
module variables it has defined, and their values: a variable defined
earlier is unbound afterwards, and using it is an error. It is rarely
needed.

\subsection{Errors and the Lisp debugger}\label{sec:lisp-errors}

\index{debugger}\index{restart}
An error that the Dylan program does not handle itself reaches SBCL's
debugger. The debugger lists the \term{restarts} available, and the Dylan
restarts in effect appear among them. For an error Dylan signals, the
loop provides ``Return to the read-eval-print loop'', which abandons the
form, prints \code{;Aborted} and gives you a new \code{?} prompt; for
every error, whatever signalled it, it also provides \code{DYLAN-REP},
``Return to the Dylan read-eval-print loop'', which does the same
without the \code{;Aborted}:

\begin{repl}
? (define x 10)
;Value: x

? (error "Oops: ~a" x)

debugger invoked on a DYLAN::DYLAN-CONDITION in thread
#<THREAD tid=690943 "main thread" RUNNING {1200030003}>:
  Error: Oops: 10

Type HELP for debugger help, or (SB-EXT:EXIT) to exit from SBCL.

restarts (invokable by number or by possibly-abbreviated name):
  0: [DYLAN-RESTART] Return to the read-eval-print loop.
  1: [DYLAN-REP    ] Return to the Dylan read-eval-print loop.
  2: [ABORT        ] Exit debugger, returning to top level.

((LAMBDA NIL :IN DYLAN::IMPLEMENTATION-SPECIFIC/INDUCE-ERROR))
0] 0
;Aborted

? x
;Value: 10
\end{repl}

\noindent
While in the debugger you can use all of SBCL's debugger commands, such
as \code{backtrace}; the frames show the Lisp the Dylan was compiled to.

A restart the program itself establishes is offered too. \code{cerror}
makes one whose description is its first argument; choosing it makes
\code{cerror} return \code{\#f} and the program carry on:

\begin{repl}
? (cerror "Use 0 instead" "bad value ~a" 7)

debugger invoked on a DYLAN::DYLAN-CONDITION ...
  Error: bad value 7
...
restarts (invokable by number or by possibly-abbreviated name):
  0: [DYLAN-RESTART] Use 0 instead
  1:                 Return to the read-eval-print loop.
  2: [DYLAN-REP    ] Return to the Dylan read-eval-print loop.
  3: [ABORT        ] Exit debugger, returning to top level.
\end{repl}

\noindent
Calling \code{(abort)} from Dylan, or signaling any \code{<abort>},
returns straight to the prompt with \code{;Aborted}, as
\dylanbook{chapter 13} describes. \code{(break "look at \textasciitilde a" 42)} enters the
debugger with the message \code{Dylan breakpoint: look at 42}; its
\code{CONTINUE} restart makes \code{break} return \code{\#f}.

\begin{notebox}
Some errors are detected by Common Lisp itself rather than by the
Dylan runtime, for example \code{(car 5)}, which reaches Lisp's
\code{car} with a non-list. Such an error arrives in the debugger as a
Lisp \code{TYPE-ERROR}, and the restarts offered are \code{DYLAN-REP}
and Lisp's own:
\begin{repl}
restarts (invokable by number or by possibly-abbreviated name):
  0: [DYLAN-REP] Return to the Dylan read-eval-print loop.
  1: [ABORT    ] Exit debugger, returning to top level.
\end{repl}
\noindent Choose \code{DYLAN-REP} to carry on at the \code{?} prompt.
(In \file{bin/dylan}, \code{ABORT} also returns to the loop. In a Lisp you
loaded the system into, it lands at the Lisp \code{*} prompt; nothing is
lost: type \code{(dylan:rep)} to go back in.) A Dylan
handler still catches these errors as instances of \code{<error>}:
\code{(handler-case (car 5) ((<error>) 'caught))} answers \code{caught}.
\end{notebox}

\subsection{Running a file at the loop}\label{sec:lisp-test-file}

\index{test@\code{dylan:test}}
The loop reads from standard input, so \cmd{bin/dylan < fact.dyl} runs
a file, printing each value as if it had been typed; at the end of the
file the program exits. The loop has no command to load a file of Dylan.
The regression test driver can do it: \code{dylan:test} reads a
file form by form, echoes each form, runs it, and prints its value, or
the message of an error, and goes on to the next form.

\begin{shell}
$ cd src
$ sbcl --eval '(load "load-dylan.lisp")' --eval '(load "full-test.lisp")' \
       --eval '(dylan:test "/path/to/fact.dyl")'
\end{shell}

\begin{repl}
(define-method fact ((n <integer>))
 (if (= n 0)
     1
     (* n (fact (- n 1)))))
;Value: fact

(print (fact 30))
265252859812191058636308480000000
;Value: 265252859812191058636308480000000

(car 3)
;Error: Lisp condition:  The value
                   3
                 is not of type
                   LIST
\end{repl}

\noindent
The definitions remain afterwards, so \code{(dylan:rep)} can follow in
the same session. Note that \file{full-test.lisp} must be loaded after
\file{load-dylan.lisp}.

\section{Compiling Dylan to Lisp files}\label{sec:lisp-compile-file}

\index{compile-file@\code{dylan:compile-file}}\index{compile-to-file@\code{dylan:compile-to-file}}
The compiler can also write its output to a file of Lisp instead of
running it, so that a program can be compiled once and loaded later
without being compiled again.

\begin{description}
\item[\code{(dylan:compile-file \meta{input} \meta{output})}] compiles
  the file of Dylan \meta{input} into the file of Lisp \meta{output}, in
  the module the input's \code{module:} header names, or
  \code{dylan-user}.
\item[\code{(dylan:compile-to-file \meta{output} \meta{form} ...)}]
  compiles Dylan forms, already read, into \meta{output}.
\end{description}

\noindent
\code{dylan:compile-file} shadows Common Lisp's function of the same
name: inside the \code{DYLAN} package, \code{compile-file} compiles
Dylan. The forms given to \code{dylan:compile-to-file} must be Dylan
forms whose symbols are in the \code{DYLAN} package, which is what the
Dylan reader \code{dylan::dylan-read} produces; Lisp forms typed at the
\code{*} prompt would have their symbols in the wrong package.

Here is a complete round trip. The program:

\begin{dylan}
;;; hello.dyl
(define-method greet ((name <string>))
  (display "Hello, ")
  (display name)
  (display "!")
  (newline))
(greet "world")
\end{dylan}

\noindent
is compiled in one Lisp session and run in another:

\begin{shell}
$ sbcl --non-interactive --eval '(load "src/load-compiler.lisp")' \
       --eval '(dylan:compile-file "hello.dyl" "hello.lisp")'
$ sbcl --non-interactive --eval '(load "src/load-runtime.lisp")' \
       --eval '(load "hello.lisp")'
Hello, world!
\end{shell}

\noindent
The output file is ordinary Common Lisp. It begins with the Dylan it
came from, as comments, and then the compiled forms:

\begin{lisp}
(in-package #:dylan)

(defvar !module/greet/ '<<the-unassigned-value>>)
...
(dylan%catch-all-conditions
 (lambda ()
   (progn
    (progn
     (cond
      ((dylan%eq? !module/greet/ '<<the-unassigned-value>>)
       (setq !module/greet/ (dylan%create-generic-function 'greet 1 nil nil)))
      ...)
     (dylan%add-method !module/greet/ ...)
     'greet)
    (dylan%call nil (dylan%free-variable-ref !module/greet/ 'greet) "world"))))
\end{lisp}

\noindent
\code{compile-to-file} is the way to compile forms you build in Lisp:

\begin{lisp}
(dylan:compile-to-file "two.lisp"
  (dylan::dylan-read (make-string-input-stream "(define x 41)"))
  (dylan::dylan-read (make-string-input-stream "(print (+ x 1))")))
\end{lisp}

\noindent
Loading \file{two.lisp} prints \code{42}. A program of several files
is compiled file by file in one session and loaded in the same order
(Section~\ref{sec:modules-lisp}).

\section{Dylan values as Lisp objects}\label{sec:lisp-representation}

\index{interactive implementation!representation}
Most Dylan values are the corresponding Lisp values, which is what makes
calling between the two languages easy.

\begin{longtable}{@{}>{\raggedright\arraybackslash}p{0.34\textwidth}>{\raggedright\arraybackslash}p{0.6\textwidth}@{}}
\toprule
Dylan & Common Lisp \\
\midrule
\endhead
integers & integers, fixnums or bignums: there is no limit on size \\
ratios & ratios \\
floats & \code{double-float}s \\
complex numbers & complex numbers; \code{(sqrt -4)} is \code{\#c(0.0 2.0)} \\
characters & characters \\
strings & Lisp strings \\
\code{<simple-object-vector>} & simple vectors \\
lists & lists; the empty list is \code{NIL} \\
\code{\#t} & \code{T} \\
\code{\#f} & a distinct object of its own, the value of \code{DYLAN::DYLAN\%FALSE} \\
symbols & symbols in the \code{DYLAN} package \\
keywords & Lisp keywords: \code{foo:} is \code{:FOO} \\
functions & Lisp functions, which take two hidden leading arguments \\
classes, instances & structures of the runtime's own \\
tables, stretchy vectors, ranges, deques & Dylan instances \\
\bottomrule
\end{longtable}

\paragraph{\code{\#f} and \code{NIL}.}\index{false@\code{\#f}}
Thomas ran on a Scheme in which false and the empty list were the same
object. The 1992 manual makes them different: \code{\#f} is the only
false value, and the empty list is true \dylanbook{chapters 3 and 8}. This
implementation follows the manual. The empty list is Lisp's \code{NIL},
because Dylan lists are Lisp lists, so \code{\#f} has to be something
else, an object created for the purpose. \code{\#t} is Lisp's \code{T}.
The consequence for Lisp code is that a Lisp ``false'' (\code{NIL}) is a
Dylan empty list and therefore \emph{true} in Dylan, and Dylan's
\code{\#f} is a non-\code{NIL} object and therefore \emph{true} in Lisp.

\paragraph{Symbols.} The Dylan reader interns each name in the
\code{DYLAN} package, folded to upper case as Lisp's reader would, so
Dylan's \code{foo} is the Lisp symbol \code{DYLAN::FOO}. (The names
\code{t} and \code{nil} are kept in lower case, so as not to be Lisp's
\code{T} and \code{NIL}.) Because the \code{DYLAN} package uses
\code{COMMON-LISP}, a Dylan symbol whose name is that of a Common Lisp
function or variable \emph{is} the Common Lisp symbol: Dylan's
\code{'expt} is \code{CL:EXPT}. The exceptions are the few names the
package shadows: \code{add-method}, \code{adjoin}, \code{compile-file},
\code{find-method}, \code{last}, \code{reduce}, \code{set-difference} and
\code{union}. The spelling a symbol was first typed in is remembered for
printing.

\section{Calling Common Lisp from Dylan}\label{sec:lisp-interop}

\index{lisp-variable}\index{lisp-procedure}
Two functions, extensions carried over from Thomas, reach into the
host Lisp.

\begin{description}
\item[\code{(lisp-variable \meta{symbol})}] answers the value of the
  Lisp global variable named by \meta{symbol}.
\item[\code{(lisp-procedure \meta{symbol})}] answers a Dylan function
  that calls the Lisp function named by \meta{symbol} (or, if the symbol
  has no function, the function that is the symbol's value). The Dylan
  function accepts any number of arguments; the Lisp function reports
  the wrong number.
\end{description}

\begin{repl}
? (lisp-variable 'most-positive-fixnum)
;Value: 4611686018427387903

? (lisp-variable 'pi)
;Value: 3.141592653589793

? ((lisp-procedure 'expt) 2 100)
;Value: 1267650600228229401496703205376

? ((lisp-procedure 'string-upcase) "hello")
;Value: "HELLO"

? (define cl-sort (lisp-procedure 'sort))
;Value: cl-sort

? (cl-sort (list 3 1 2) '<)
;Value: (1 2 3)
\end{repl}

\noindent
Functions of your own written in Lisp can be called the same way. Define
them in the \code{DYLAN} package, so that the Dylan symbol names them:

\begin{shell}
$ sbcl --eval '(load "src/load-dylan.lisp")' \
       --eval '(defun dylan::square (x) (* x x))' --eval '(dylan:rep)'
\end{shell}

\begin{repl}
? (define square (lisp-procedure 'square))
;Value: square

? (square 12)
;Value: 144

? (square 1 2)

debugger invoked on a DYLAN::DYLAN-CONDITION ...
  Error: invalid number of arguments: 2
\end{repl}

\noindent
Keep these rules in mind:

\begin{itemize}
\item \textbf{Lisp's false is Dylan's true.} A Lisp predicate answers
  \code{NIL} for false, which Dylan sees as the empty list, a true
  value:
\begin{dylan}
(if ((lisp-procedure 'evenp) 3) 'yes 'no)    ; => yes
\end{dylan}
  Test the result against the empty list instead:
\begin{dylan}
(define lisp-true? (method (x) (not (id? x '()))))
(lisp-true? ((lisp-procedure 'evenp) 3))      ; => #f
\end{dylan}
  Likewise, do not pass \code{\#f} to a Lisp function that expects
  \code{NIL}; pass \code{'()}.
\item \textbf{Dylan functions cannot be passed to Lisp.} A Dylan
  function, including one made by \code{lisp-procedure}, expects two
  hidden arguments that Lisp does not supply, so handing one to a Lisp
  function such as \code{mapcar} fails with ``invalid number of
  arguments''. Pass a symbol naming a Lisp function, as with \code{'<}
  above.
\item \textbf{Lisp values print as Dylan sees them.} A Lisp value that
  has no Dylan meaning, such as a package or a hash table, prints in
  Lisp's notation.
\end{itemize}

\section{How it is made fast}\label{sec:lisp-performance}

\index{interactive implementation!performance}
Thomas's authors did not aim at speed, and a direct translation of it
was slow. The Common Lisp implementation has since been rebuilt around a
few techniques, all of which keep the language's full dynamism.

\begin{itemize}
\item \textbf{Dispatch caches.} Each generic function caches, for each
  combination of argument classes it has seen, the sorted list of
  applicable methods, along with the keyword checks those methods need.
  Adding or removing a method, or changing a class's superclasses,
  invalidates the cache. A generic function of one, two or three
  required arguments calls the cached method directly without building
  an argument list, and calls with several applicable methods
  (\code{next-method} chains) and with keyword arguments go through the
  cache too. Singleton specializers are found through a hash table.
\item \textbf{Direct calls.} A call compiles to a direct call of the
  function, without wrapping each argument in a closure as Thomas did. A
  method whose specializers are all \code{<object>} is called as its
  compiled body, without an argument-checking wrapper.
\item \textbf{Guarded fast paths.} Calls of \code{+}, \code{-},
  \code{*}, \code{<}, \code{>}, \code{<=}, \code{>=}, \code{=},
  \code{modulo}, \code{remainder}, \code{element} and
  \code{(setter element)} compile inline when the arguments are Lisp
  numbers or characters, or a simple vector or string and an index.
  \code{for-each}, \code{map}, \code{reduce} and their kin walk lists,
  simple vectors and ranges directly. Each fast path is guarded by a
  flag that is true only while no method a program has added could apply
  to those arguments; add a method to \code{binary+} that could apply to
  integers, and the inline path for \code{+} turns itself off, so the
  language's meaning never changes.
\item \textbf{Direct slot access.} Slots the runtime's own classes keep
  privately are read straight from the instance, and each class caches
  its slots' required init keywords for \code{make}.
\end{itemize}

\noindent
\file{generic.lisp} holds the dispatcher and its caches;
\file{runtime-fast-paths.lisp} turns the fast paths on once every
predefined method exists. \file{bench/run.lisp} measures the
implementation against the same programs written in plain Common Lisp;
Chapter~\ref{ch:performance} gives the results.

\section{The book's examples and the regression test}\label{sec:lisp-examples}

\index{regression test}\index{dylan-examples.dyl@\file{dylan-examples.dyl}}
The file \file{dylan-examples.dyl} at the top of the repository holds the
examples of the 1992 manual, extracted and edited into runnable form.
\code{dylan:test-examples} runs it with \code{dylan:test}
(Section~\ref{sec:lisp-test-file}):

\begin{shell}
$ cd src
$ sbcl --eval '(load "load-dylan.lisp")' \
       --eval '(load "full-test.lisp")' \
       --eval '(dylan:test-examples)'
\end{shell}

\noindent
The transcript it prints is compared with \file{src/reference/MIT.output},
the transcript Thomas produced from the same file on MIT Scheme in 1992
(the transcripts from Gambit and Scheme->C are kept beside it). The two
agree on every value except a handful of places where the manual and
Thomas disagree and the manual is followed: \code{\#f} and \code{()} are
different objects, floats print as \code{9.0} rather than \code{9.},
\code{(sqrt 25)} is \code{5.0}, \code{make} given a keyword that
nothing accepts is an error (the book's \code{<displaced-view>} example
does this, so it and the forms after it fail), and \code{map-into}
passes only the source elements. \file{src/README} and
Appendix~\ref{app:testing} list them.

\section{Differences from the compiler to C}\label{sec:lisp-vs-dylanc}

\index{interactive implementation!versus compiler to C}
The two implementations accept the same language, but each goes beyond
the other in places.

\paragraph{What only the interactive implementation does.} It is fully
dynamic, as the 1992 manual allows:

\begin{itemize}
\item \code{add-method}, \code{remove-method}, \code{add-slot},
  \code{remove-slot} and \code{freeze-methods} work. A program can add a
  slot to a class that already has instances:
\begin{repl}
? (define-class <sq> (<object>) (side init-keyword: side:))
;Value: <sq>

? (define s (make <sq> side: 2))
;Value: s

? (define-generic-function color (object))
;Value: color

? (add-slot <sq> getter: color init-value: 'red)
;Value: {slot color}

? (color s)
;Value: red
\end{repl}
\item New code can be compiled and run at any time, and any definition
  can be replaced (Section~\ref{sec:lisp-rep}).
\item Integers have no limit on size, and complex numbers exist:
  \code{(* 99999999999 99999999999)} is \code{9999999999800000000001},
  and \code{(sqrt -4)} is \code{\#c(0.0 2.0)}. The compiler to C has
  63-bit integers, where overflow is an error, and no complex numbers.
\item \code{lisp-variable} and \code{lisp-procedure} reach the host
  Lisp.
\end{itemize}

\paragraph{What only the compiler to C does.} Its extensions are
absent here:

\begin{itemize}
\item no threads, locks, condition variables, atomic integers, thread
  variables or \code{dynamic-bind};
\item no macros: \code{define-macro} and \code{define-for-macros} are
  not special forms here;
\item no \code{current-milliseconds}.
\end{itemize}

\noindent
Appendix~\ref{app:library} marks every name that only one
implementation provides.

\paragraph{What behaves differently.}

\begin{itemize}
\item Modules must be declared before they are used, in order; the
  compiler to C sees the whole program first (Chapter~\ref{ch:modules}).
\item Errors are reported when the code runs; the compiler to C reports
  many of them, such as an assignment to a read-only variable, when it
  compiles.
\end{itemize}

\section{Limitations}\label{sec:lisp-limitations}

\begin{itemize}
\item There is no way to load a file of Dylan source at the
  read-eval-print loop itself; use \code{dylan:test} or compile the file
  with \code{dylan:compile-file} and load the result.
\item Thomas's MIT Scheme restart example, written against MIT Scheme's
  condition system, has no counterpart; it would have to be rewritten
  against Common Lisp's.
\end{itemize}
